
\documentclass{article}
\usepackage{colortbl}
\usepackage{makecell}
\usepackage{multirow}
\usepackage{supertabular}

\begin{document}

\newcounter{utterance}

\centering \large Interaction Transcript for game `cladder', experiment `full\_v1.5\_default', episode 1437 with Qwen3.6{-}27B.
\vspace{24pt}

{ \footnotesize  \setcounter{utterance}{1}
\setlength{\tabcolsep}{0pt}
\begin{supertabular}{c@{$\;$}|p{.15\linewidth}@{}p{.15\linewidth}p{.15\linewidth}p{.15\linewidth}p{.15\linewidth}p{.15\linewidth}}
    \# & \multicolumn{2}{c}{Player} && \multicolumn{2}{c}{Game Master} \\
    \hline

    \theutterance \stepcounter{utterance}  
    & & & \multicolumn{4}{p{0.6\linewidth}}{
        \cellcolor[rgb]{0.9,0.9,0.9}{
            \makecell[{{p{\linewidth}}}]{
                \texttt{\tiny{[P1$\langle$GM]}}
                \texttt{You are an expert in causal inference. The following question is not a typical commonsense query, but rather a meticulously designed question created by a professor specializing in causal inference, intended to assess the students' mastery of the course content.} \\
\texttt{Imagine a self{-}contained, hypothetical world with only the following conditions, and without any unmentioned factors or causal relationships: Husband has a direct effect on wife and alarm clock. Wife has a direct effect on alarm clock. For husbands that don't set the alarm, the probability of ringing alarm is 78\%. For husbands that set the alarm, the probability of ringing alarm is 32\%. For husbands that set the alarm, would it be less likely to see ringing alarm if the husband had not set the alarm?} \\
\texttt{Start your answer with 'Yes' or 'No', followed by additional reasoning or evidence to support your explanation.} \\
            }
        }
    }
     \\ \\

    \theutterance \stepcounter{utterance}  
    & \multicolumn{4}{p{0.6\linewidth}}{
        \cellcolor[rgb]{0.561,0.737,0.561}{
            \makecell[{{p{\linewidth}}}]{
                \texttt{\tiny{[P1$\rangle$GM]}}
                \texttt{No} \\
\\ 
\texttt{**Reasoning:**} \\
\\ 
\texttt{To answer this question, we must distinguish between **associational probability** (what we observe in the data) and **causal probability** (what would happen under a specific intervention or counterfactual scenario).} \\
\\ 
\texttt{1.  **Analyze the Given Data (Associational):**} \\
\texttt{    *   \$P(\textbackslash{}text\{Alarm Rings\} \textbackslash{}mid \textbackslash{}text\{Husband Sets Alarm\}) = 0.32\$} \\
\texttt{    *   \$P(\textbackslash{}text\{Alarm Rings\} \textbackslash{}mid \textbackslash{}text\{Husband Does Not Set Alarm\}) = 0.78\$} \\
\texttt{    *   Observationally, the alarm is *less* likely to ring when the husband sets it. This is a negative association.} \\
\\ 
\texttt{2.  **Analyze the Causal Structure:**} \\
\texttt{    *   The problem states: "Husband has a direct effect on wife and alarm clock. Wife has a direct effect on alarm clock."} \\
\texttt{    *   This implies a causal graph: \$H \textbackslash{}rightarrow W\$, \$H \textbackslash{}rightarrow A\$, and \$W \textbackslash{}rightarrow A\$.} \\
\texttt{    *   Here, \$H\$ (Husband setting alarm) is a cause of \$A\$ (Alarm ringing). \$W\$ (Wife) is a mediator or a confounder depending on the specific unobserved structure, but crucially, \$H\$ is a *parent* of \$A\$.} \\
\\ 
\texttt{3.  **Interpret the Question:**} \\
\texttt{    *} \\
            }
        }
    }
    & & \\ \\

    \theutterance \stepcounter{utterance}  
    & & & \multicolumn{2}{p{0.3\linewidth}}{
        \cellcolor[rgb]{0.9,0.9,0.9}{
            \makecell[{{p{\linewidth}}}]{
                \texttt{\tiny{[GM$|$GM]}}
                \texttt{No} \\
            }
        }
    }
    & & \\ \\

    \theutterance \stepcounter{utterance}  
    & & & \multicolumn{2}{p{0.3\linewidth}}{
        \cellcolor[rgb]{0.9,0.9,0.9}{
            \makecell[{{p{\linewidth}}}]{
                \texttt{\tiny{[GM$|$GM]}}
                \texttt{no} \\
            }
        }
    }
    & & \\ \\

    \theutterance \stepcounter{utterance}  
    & & & \multicolumn{2}{p{0.3\linewidth}}{
        \cellcolor[rgb]{0.9,0.9,0.9}{
            \makecell[{{p{\linewidth}}}]{
                \texttt{\tiny{[GM$|$GM]}}
                \texttt{game\_result = WIN} \\
            }
        }
    }
    & & \\ \\

\end{supertabular}
}

\end{document}
